\BookPart{Encoding and molecular algorithms}{algorithms}{%
With the toolkit in hand we can design computations deliberately. Part~III
asks how to choose sequences that mean what we intend and only that, how to
organise filtering algorithms so that an accepted molecule is provably a
solution, and how far the generate-filter-verify paradigm stretches: to Boolean
satisfiability and to rewritable molecular memories. Throughout, a
physical population is a multiset of finite copies, and every filter is a noisy
observation rather than a perfect predicate.

\textbf{Chapters in this part}
\begin{itemize}\setlength{\itemsep}{2pt}
\item \textbf{11\quad Sequence design and graph encoding.} How do we choose codewords that pair only with their intended partners?
\item \textbf{12\quad Generate--filter--verify algorithms.} When is a molecular filtering algorithm sound and complete, and what do losses do to it?
\item \textbf{13\quad SAT and combinatorial constructions.} How can a tube decide a Boolean formula, and where does the exponential cost reappear?
\item \textbf{14\quad Sticker systems and molecular memory.} How can bits be written into and read from a population of memory strands?
\end{itemize}}
